\section{Conclusion}

Alfred Snider first recognized debate as a game and championed its transformative power \parencite{snider1982gaming,snider1984gaming}, yet none who followed applied game-design know-how to the question his argument makes inevitable: how can we make it a better game? This thesis takes one step into that question.

This thesis pursues that ``how can'' question in three movements: establishing that the conceptual tools apply (Sections \ref{sec:background}--\ref{sec:theory}: debate is a transformative role-playing game, and game design can be brought to bear on it); carrying a concrete design through (Sections \ref{sec:methods}--\ref{sec:design}: the Inside Out format and its iterative development); and reading what that design evoked in play (Results and Discussion: a clue-by-clue account and its theory-facing extensions).

Its contribution to the broader field of transformative game design is at the level of proof of concept: it presents a design and draws from the limited playtest data an account of how that design shapes the experience of those who play it, and where it shows promise.

A word on how the work was actually made. The design did not deduce itself from theory; it emerged through iteration and intuition, and the theoretical analysis came alongside and after---illuminating why the moves that felt right worked, and where they might be pushed further. Under a hypothesis-testing ideal this order would be awkward to confess. Under Research through Design it is simply the method: design leads, theory follows and travels with it, and a design that holds up under theoretical scrutiny has earned something rather than faked it. What the scholarly frame added was not a deductive justification after the fact but a way of making the design's reasoning legible and transferable. More playtest cycles, rather than fewer, would have sharpened the format further; the binding constraint was the difficulty of recruiting adult playtesters, not a limit of the approach.

Inside Out is a first---and not strictly necessary---step in a much longer process of reform, whose later stages involve refinement, community adoption, governance design, and the slow work of shifting collective habits. What the thesis carries forward into that process is threefold: a working design and the documented reasoning behind it, ready for the next cycle; the momentum to pursue the stages of reform the design alone cannot deliver; and a scholarly articulation that makes the project legible to others who might join it---those at the uncommon intersection of debate heritage, game-design practice, and scholarly ambition.

The longer aim, beyond this thesis, is to build something financially sustainable that orchestrates and supports debate-centric educational gameplay; the first steps toward it were named at the end of the discussion. Alfred ``Tuna'' Snider built a life of meaningful teaching, reaching hundreds of young people, on the conviction that debate is a game worth taking seriously as one. That precedent is reason enough to keep going.